{\renewcommand{\baselinestretch}{1.2}\normalsize
大型視覺語言模型（Vision-Language Models，VLMs）在多模態任務上表現卓越，但高度依賴雲端基礎設施與高效能 GPU，難以直接部署於資源受限的邊緣環境。本論文針對記憶體低於 800 MB、純 CPU 推理、延遲低於 10 毫秒及任務關鍵場景零幻覺容忍等嚴苛約束，開發適用於裝置端部署的輕量視覺語言模型。

本論文提出兩套系統。\textbf{BitMar} 為 1400 萬參數、採用 1.58 位元三元量化的多模態模型，整合跨模態融合管線與外部情節記憶矩陣（512 槽位，128 維）。啟用情節記憶後，推理吞吐量提升約 7.5 倍，能耗降低 79\%，並在實體追蹤、組合推理及視覺問答任務上獲得 3 至 4 個百分點的準確率提升，驗證了極限量化、跨模態融合與情節記憶在同一架構中共存的可行性。\textbf{EmberVLM} 為本論文主要貢獻，總計 1.64 億參數（4000 萬可訓練），採用四階段課程式訓練策略。其情節記憶模組採用高斯距離定址與新穎性觸發動態寫入機制，每步推理僅增加 1.2 MB 記憶體佔用及不足 3\% 的浮點運算量，可將機器人選擇任務之巨集 F1 分數從 0.035 提升至 0.297。其視覺缺失精煉器（VA Refiner）為免訓練的幻覺抑制機制，透過雙視角邏輯差異分析與解碼器中間層神經元激活監控，無需額外前向傳遞即可抑制視覺上無依據的輸出。

EmberVLM 在三項任務關鍵基準上取得最佳成績：Top-N 機器人選擇（巨集 F1 = 0.42）、緊急視覺辨識（準確率 0.64）及地理災害理解（準確率 0.72）。完整推理堆疊在純 CPU 裝置上僅需 5.22 毫秒，記憶體佔用低於 800 MB。VA 精煉器將對抗性幻覺率從 29\% 降至 22\%。全訓練週期碳排放實測為 25.3 kg CO$_2$eq，展示了任務關鍵邊緣人工智慧的永續訓練路徑。

\chinesekeywords{邊緣人工智慧、視覺語言模型、情節記憶、幻覺抑制、模型量化、資源受限部署}
}
